\section{Results}
\label{sec:results}

\subsection{Evolution Converges Rapidly}
\label{sec:convergence}

\Tabref{tab:convergence} shows the evolution trajectory.
The population reaches ceiling fitness (1.000) by generation 2, up from an average of 0.647 at initialization.
Notably, one seed prompt (Prompt~4) achieves perfect fitness even at generation 0 due to a structured GOAL/RULE/PRIORITY format, suggesting that structured formatting itself is a powerful goal-persistence feature.
A single anomalous prompt (Prompt~3) scores 0.000 despite seeming reasonable, likely because the model interpreted its declarative phrasing as a description rather than an instruction.

\begin{table}[t]
    \centering
    \begin{tabular}{@{}lccc@{}}
        \toprule
        Generation & Best Fitness & Avg Fitness & Avg Length (chars) \\
        \midrule
        0   & 1.000 & 0.647 & 148 \\
        1   & 1.000 & 0.948 & $\sim$400 \\
        2   & 1.000 & {\bf 1.000} & $\sim$800 \\
        15  & 1.000 & {\bf 1.000} & 1{,}609 \\
        \bottomrule
    \end{tabular}
    \caption{Evolution convergence over generations.
    The population reaches ceiling fitness by generation~2, while prompt length continues to grow through generation~15 as evolution adds redundancy.}
    \label{tab:convergence}
\end{table}

After convergence, evolution continues to increase prompt length (from $\sim$800 chars at generation~2 to 1{,}609 at generation~15) without improving fitness, suggesting that the algorithm adds defensive redundancy even when it is not necessary for the current task.

\subsection{Main Results: Ceiling Effect}
\label{sec:main_results}

\Tabref{tab:main_results} presents results on the held-out test set.
Both \strongelicit and all evolved prompts achieve perfect scores across all metrics.
The \zeroshot baseline performs poorly (fitness 0.096), while \simple achieves moderate performance (0.836).

\begin{table}[t]
    \centering
    \begin{tabular}{@{}lccc@{}}
        \toprule
        Prompt & Fitness & Persistence & Return Rate \\
        \midrule
        \zeroshot       & 0.096 & 0.115 & 0.067 \\
        \simple         & 0.836 & 0.860 & 0.800 \\
        \strongelicit   & {\bf 1.000} & {\bf 1.000} & {\bf 1.000} \\
        \midrule
        Evolved \#1     & {\bf 1.000} & {\bf 1.000} & {\bf 1.000} \\
        Evolved \#2     & {\bf 1.000} & {\bf 1.000} & {\bf 1.000} \\
        Evolved \#3     & {\bf 1.000} & {\bf 1.000} & {\bf 1.000} \\
        \bottomrule
    \end{tabular}
    \caption{Held-out test results (5 unseen scenarios, 3 runs each).
    Both \strongelicit and all evolved prompts achieve perfect scores.
    Best results in \textbf{bold}.}
    \label{tab:main_results}
\end{table}

The prompt quality hierarchy is stark: \zeroshot (0.096) $\ll$ \simple (0.836) $\ll$ \strongelicit (1.0) $=$ Evolved (1.0).
The 8.4$\times$ improvement from \zeroshot to \simple demonstrates that instruction quality matters enormously, while the gap from \simple to \strongelicit shows the value of emphatic goal statements and explicit distraction negation.

\subsection{Statistical Analysis}
\label{sec:stats}

\Tabref{tab:stats} summarizes our hypothesis tests.
The comparison between evolved and \strongelicit prompts yields identical zeros across all metrics, making statistical testing inapplicable ($d = 0.000$).
Evolved prompts significantly outperform \zeroshot ($W = 105.0$, $p = 0.0003$, $d = 5.291$) but the comparison with \simple does not survive Bonferroni correction ($W = 6.0$, $p = 0.042 > 0.017$).

\begin{table}[t]
    \centering
    \resizebox{\textwidth}{!}{%
    \begin{tabular}{@{}lcccc@{}}
        \toprule
        Comparison & Wilcoxon $W$ & $p$-value & Significant? & Cohen's $d$ \\
        \midrule
        Evolved vs.\ \zeroshot       & 105.0 & 0.0003 & Yes ($p < 0.017$) & 5.291 \\
        Evolved vs.\ \simple         & 6.0   & 0.042  & No ($p > 0.017$)  & 0.707 \\
        Evolved vs.\ \strongelicit   & ---   & ---    & All zeros          & 0.000 \\
        \bottomrule
    \end{tabular}%
    }
    \caption{Hypothesis test results with Bonferroni-corrected $\alpha = 0.017$.
    The evolved vs.\ \strongelicit comparison produces identical scores, yielding zero effect size.}
    \label{tab:stats}
\end{table}

\para{Hypothesis outcomes.}
H1 (evolved $>$ \strongelicit on persistence): \textbf{not supported}, $d = 0.000$.
H2 (higher return-to-goal rate): \textbf{not supported}, both achieve 100\%.
H3 (identifiable linguistic features): \textbf{supported}, as we detail in \secref{sec:features}.

\subsection{Feature Analysis: What Evolution Discovers}
\label{sec:features}

Although evolved and human prompts achieve identical performance, they differ qualitatively.
\Tabref{tab:features} compares 8 goal-persistence features across prompt types.

\begin{table}[t]
    \centering
    \resizebox{\textwidth}{!}{%
    \begin{tabular}{@{}lcccc@{}}
        \toprule
        Feature & \zeroshot & \simple & \strongelicit & Evolved \\
        \midrule
        Explicit priority         & --- & \checkmark & \checkmark & \checkmark \\
        Negation/prohibition      & --- & ---        & \checkmark & \checkmark \\
        Distraction awareness     & --- & ---        & \checkmark & \checkmark \\
        Identity framing          & --- & \checkmark & ---        & \checkmark \\
        Output format spec.       & --- & \checkmark & \checkmark & \checkmark \\
        Emotional resistance      & --- & ---        & \checkmark & \checkmark \\
        Recovery language         & --- & ---        & ---        & {\bf \checkmark} \\
        Distraction enumeration   & --- & ---        & \checkmark & \checkmark \\
        \midrule
        Features present          & 1/8 & 3/8       & 6/8        & {\bf 8/8} \\
        \bottomrule
    \end{tabular}%
    }
    \caption{Goal-persistence features across prompt types.
    Evolved prompts contain all 8 features, including \emph{recovery language} (unique to evolved prompts).}
    \label{tab:features}
\end{table}

\para{Recovery language.}
The most notable evolved feature is explicit recovery instructions: ``If interrupted, confused, or sidetracked---even instantaneously or momentarily---immediately and silently resume counting at the exact next integer with zero delay.''
This instruction is absent from all human-designed baselines, including \strongelicit.

\para{Structural convergence.}
All 10 final evolved prompts converge to a GOAL/RULE/PRIORITY/MANDATE/STRATEGY section organization.
This homogenization suggests a single attractor in the fitness landscape for this task.

\para{Redundancy as strategy.}
Evolved prompts repeat instructions in multiple phrasings (\eg ``block, ignore, discard, and suppress'' vs.\ simply ``ignore'').
At 1{,}609 characters on average, evolved prompts are approximately 4$\times$ longer than \strongelicit (378 characters).
This redundancy may serve as implicit self-reinforcement during generation, priming the model to maintain focus through repeated exposure to the goal.

\subsection{Training Set Consistency}
\label{sec:training}

Results on the training set mirror the held-out test: \zeroshot achieves 0.036 fitness, \simple achieves 0.461, and both \strongelicit and evolved prompts achieve 1.000.
The consistency between training and test performance confirms that the ceiling effect is not an artifact of overfitting.
